% Options for packages loaded elsewhere
% Options for packages loaded elsewhere
\PassOptionsToPackage{unicode}{hyperref}
\PassOptionsToPackage{hyphens}{url}
\PassOptionsToPackage{dvipsnames,svgnames,x11names}{xcolor}
%
\documentclass[
  korean,
]{article}
\usepackage{xcolor}
\usepackage{amsmath,amssymb}
\setcounter{secnumdepth}{5}
\usepackage{iftex}
\ifPDFTeX
  \usepackage[T1]{fontenc}
  \usepackage[utf8]{inputenc}
  \usepackage{textcomp} % provide euro and other symbols
\else % if luatex or xetex
  \usepackage{unicode-math} % this also loads fontspec
  \defaultfontfeatures{Scale=MatchLowercase}
  \defaultfontfeatures[\rmfamily]{Ligatures=TeX,Scale=1}
\fi
\usepackage{lmodern}
\ifPDFTeX\else
  % xetex/luatex font selection
\fi
% Use upquote if available, for straight quotes in verbatim environments
\IfFileExists{upquote.sty}{\usepackage{upquote}}{}
\IfFileExists{microtype.sty}{% use microtype if available
  \usepackage[]{microtype}
  \UseMicrotypeSet[protrusion]{basicmath} % disable protrusion for tt fonts
}{}
\makeatletter
\@ifundefined{KOMAClassName}{% if non-KOMA class
  \IfFileExists{parskip.sty}{%
    \usepackage{parskip}
  }{% else
    \setlength{\parindent}{0pt}
    \setlength{\parskip}{6pt plus 2pt minus 1pt}}
}{% if KOMA class
  \KOMAoptions{parskip=half}}
\makeatother
% Make \paragraph and \subparagraph free-standing
\makeatletter
\ifx\paragraph\undefined\else
  \let\oldparagraph\paragraph
  \renewcommand{\paragraph}{
    \@ifstar
      \xxxParagraphStar
      \xxxParagraphNoStar
  }
  \newcommand{\xxxParagraphStar}[1]{\oldparagraph*{#1}\mbox{}}
  \newcommand{\xxxParagraphNoStar}[1]{\oldparagraph{#1}\mbox{}}
\fi
\ifx\subparagraph\undefined\else
  \let\oldsubparagraph\subparagraph
  \renewcommand{\subparagraph}{
    \@ifstar
      \xxxSubParagraphStar
      \xxxSubParagraphNoStar
  }
  \newcommand{\xxxSubParagraphStar}[1]{\oldsubparagraph*{#1}\mbox{}}
  \newcommand{\xxxSubParagraphNoStar}[1]{\oldsubparagraph{#1}\mbox{}}
\fi
\makeatother


\usepackage{longtable,booktabs,array}
\newcounter{none} % for unnumbered tables
\usepackage{calc} % for calculating minipage widths
% Correct order of tables after \paragraph or \subparagraph
\usepackage{etoolbox}
\makeatletter
\patchcmd\longtable{\par}{\if@noskipsec\mbox{}\fi\par}{}{}
\makeatother
% Allow footnotes in longtable head/foot
\IfFileExists{footnotehyper.sty}{\usepackage{footnotehyper}}{\usepackage{footnote}}
\makesavenoteenv{longtable}
\usepackage{graphicx}
\makeatletter
\newsavebox\pandoc@box
\newcommand*\pandocbounded[1]{% scales image to fit in text height/width
  \sbox\pandoc@box{#1}%
  \Gscale@div\@tempa{\textheight}{\dimexpr\ht\pandoc@box+\dp\pandoc@box\relax}%
  \Gscale@div\@tempb{\linewidth}{\wd\pandoc@box}%
  \ifdim\@tempb\p@<\@tempa\p@\let\@tempa\@tempb\fi% select the smaller of both
  \ifdim\@tempa\p@<\p@\scalebox{\@tempa}{\usebox\pandoc@box}%
  \else\usebox{\pandoc@box}%
  \fi%
}
% Set default figure placement to htbp
\def\fps@figure{htbp}
\makeatother

\ifLuaTeX
  \usepackage{luacolor}
  \usepackage[soul]{lua-ul}
\else
  \usepackage{soul}
\fi


\ifLuaTeX
\usepackage[bidi=basic,shorthands=off,provide=*]{babel}
\else
\usepackage[bidi=default,shorthands=off,provide=*]{babel}
\fi
\ifLuaTeX
  \usepackage{selnolig} % disable illegal ligatures
\fi


\setlength{\emergencystretch}{3em} % prevent overfull lines

\providecommand{\tightlist}{%
  \setlength{\itemsep}{0pt}\setlength{\parskip}{0pt}}



 


\makeatletter
\@ifpackageloaded{tcolorbox}{}{\usepackage[skins,breakable]{tcolorbox}}
\@ifpackageloaded{fontawesome5}{}{\usepackage{fontawesome5}}
\definecolor{quarto-callout-color}{HTML}{909090}
\definecolor{quarto-callout-note-color}{HTML}{0758E5}
\definecolor{quarto-callout-important-color}{HTML}{CC1914}
\definecolor{quarto-callout-warning-color}{HTML}{EB9113}
\definecolor{quarto-callout-tip-color}{HTML}{00A047}
\definecolor{quarto-callout-caution-color}{HTML}{FC5300}
\definecolor{quarto-callout-color-frame}{HTML}{acacac}
\definecolor{quarto-callout-note-color-frame}{HTML}{4582ec}
\definecolor{quarto-callout-important-color-frame}{HTML}{d9534f}
\definecolor{quarto-callout-warning-color-frame}{HTML}{f0ad4e}
\definecolor{quarto-callout-tip-color-frame}{HTML}{02b875}
\definecolor{quarto-callout-caution-color-frame}{HTML}{fd7e14}
\makeatother
\makeatletter
\@ifpackageloaded{caption}{}{\usepackage{caption}}
\AtBeginDocument{%
\ifdefined\contentsname
  \renewcommand*\contentsname{목차}
\else
  \newcommand\contentsname{목차}
\fi
\ifdefined\listfigurename
  \renewcommand*\listfigurename{그림 목록}
\else
  \newcommand\listfigurename{그림 목록}
\fi
\ifdefined\listtablename
  \renewcommand*\listtablename{표 목록}
\else
  \newcommand\listtablename{표 목록}
\fi
\ifdefined\figurename
  \renewcommand*\figurename{그림}
\else
  \newcommand\figurename{그림}
\fi
\ifdefined\tablename
  \renewcommand*\tablename{표}
\else
  \newcommand\tablename{표}
\fi
}
\@ifpackageloaded{float}{}{\usepackage{float}}
\floatstyle{ruled}
\@ifundefined{c@chapter}{\newfloat{codelisting}{h}{lop}}{\newfloat{codelisting}{h}{lop}[chapter]}
\floatname{codelisting}{목록}
\newcommand*\listoflistings{\listof{codelisting}{코드 목록}}
\makeatother
\makeatletter
\makeatother
\makeatletter
\@ifpackageloaded{caption}{}{\usepackage{caption}}
\@ifpackageloaded{subcaption}{}{\usepackage{subcaption}}
\makeatother
\usepackage{bookmark}
\IfFileExists{xurl.sty}{\usepackage{xurl}}{} % add URL line breaks if available
\urlstyle{same}
\hypersetup{
  pdftitle={Insertion Sort를 Incremental Algorithm 관점에서 이해하기},
  pdflang={ko},
  colorlinks=true,
  linkcolor={blue},
  filecolor={Maroon},
  citecolor={Blue},
  urlcolor={Blue},
  pdfcreator={LaTeX via pandoc}}


\title{Insertion Sort를 Incremental Algorithm 관점에서 이해하기}
\author{}
\date{}
\begin{document}
\maketitle


\section{서론}\label{uxc11cuxb860}

Insertion Sort는 가장 기본적인 정렬 알고리즘 중 하나이며,\\
\textbf{incremental algorithm(점진적 알고리즘)} 의 대표적인 예시로 자주
소개된다.

Insertion Sort를 incremental algorithm이라고 부르는 이유는,\\
각 단계에서 이미 해결된 작은 문제의 해를 바탕으로 다음 단계의 해를
만들기 때문이다.\\

각 반복 단계에서 알고리즘은 다음과 같이 생각한다.\\

현재 A{[}1..j-1{]}은 이미 정렬되어 있다.\\
이제 새로운 원소 A{[}j{]}를 그 정렬된 구간 안의 적절한 위치에 넣는다.\\
그러면 A{[}1..j{]}도 정렬된다. 즉, 정렬은 다음과 같은 순서로 점점
확장된다.\\

길이 1의 정렬된 배열\\
길이 2의 정렬된 배열\\
길이 3의 정렬된 배열\\
\ldots{}\\
길이 n의 정렬된 배열\\

이처럼 정렬된 부분배열을 한 칸씩 확장한다 는 점이 바로
\ul{\textbf{incremental algorithm}}의 본질이다.\\

incremental algorithm의 핵심 아이디어는 다음과 같다.

\begin{tcolorbox}[enhanced jigsaw, arc=.35mm, bottomrule=.15mm, bottomtitle=1mm, breakable, colback=white, colbacktitle=quarto-callout-note-color!10!white, colframe=quarto-callout-note-color-frame, coltitle=black, left=2mm, leftrule=.75mm, opacityback=0, opacitybacktitle=0.6, rightrule=.15mm, title=\textcolor{quarto-callout-note-color}{\faInfo}\hspace{0.5em}{Incremental Algorithm의 핵심 아이디어}, titlerule=0mm, toprule=.15mm, toptitle=1mm]

전체 문제를 한 번에 해결하는 대신,\\
\textbf{이미 해결된 더 작은 부분문제의 해를 바탕으로 새로운 원소를
하나씩 추가하여 전체 해를 확장해 나가는 방식}이다.

\end{tcolorbox}

Insertion Sort는 바로 이런 방식으로 동작한다.\\
즉, 이미 정렬된 부분배열에 새로운 원소를 적절한 위치에 삽입하면서\\
정렬된 구간을 점점 확장해 나간다.

\begin{longtable}[]{@{}llll@{}}
\caption{Insertion Sort 동작 원리}\tabularnewline
\toprule\noalign{}
\textbf{단계} & \textbf{key} & \textbf{삽입 전 정렬 구간} & \textbf{삽입
후 배열} \\
\midrule\noalign{}
\endfirsthead
\toprule\noalign{}
\textbf{단계} & \textbf{key} & \textbf{삽입 전 정렬 구간} & \textbf{삽입
후 배열} \\
\midrule\noalign{}
\endhead
\bottomrule\noalign{}
\endlastfoot
시작 & - & \texttt{{[}5{]}} & \texttt{{[}5,\ 2,\ 4,\ 6,\ 1,\ 3{]}} \\
1 & 2 & \texttt{{[}5{]}} & \texttt{{[}2,\ 5,\ 4,\ 6,\ 1,\ 3{]}} \\
2 & 4 & \texttt{{[}2,\ 5{]}} & \texttt{{[}2,\ 4,\ 5,\ 6,\ 1,\ 3{]}} \\
3 & 6 & \texttt{{[}2,\ 4,\ 5{]}} &
\texttt{{[}2,\ 4,\ 5,\ 6,\ 1,\ 3{]}} \\
4 & 1 & \texttt{{[}2,\ 4,\ 5,\ 6{]}} &
\texttt{{[}1,\ 2,\ 4,\ 5,\ 6,\ 3{]}} \\
5 & 3 & \texttt{{[}1,\ 2,\ 4,\ 5,\ 6{]}} &
\texttt{{[}1,\ 2,\ 3,\ 4,\ 5,\ 6{]}} \\
\end{longtable}

\section{Loop Invariant로 본
정당성}\label{loop-invariantuxb85c-uxbcf8-uxc815uxb2f9uxc131}

Insertion Sort의 correctness를 설명할 때 가장 표준적인 도구는
\textbf{loop invariant} 이다.

\begin{tcolorbox}[enhanced jigsaw, arc=.35mm, bottomrule=.15mm, bottomtitle=1mm, breakable, colback=white, colbacktitle=quarto-callout-important-color!10!white, colframe=quarto-callout-important-color-frame, coltitle=black, left=2mm, leftrule=.75mm, opacityback=0, opacitybacktitle=0.6, rightrule=.15mm, title=\textcolor{quarto-callout-important-color}{\faExclamation}\hspace{0.5em}{Loop Invariant}, titlerule=0mm, toprule=.15mm, toptitle=1mm]

\texttt{for} 반복문의 각 반복이 시작될 때마다, 부분배열
\texttt{A{[}1..j-1{]}}은 원래 그 원소들로 이루어진 \textbf{정렬된
배열}이다.\\

\end{tcolorbox}

이 invariant는 incremental algorithm의 철학과 정확히 일치한다. 즉, 이전
단계까지의 해가 올바르다고 가정하고, 그 위에 새로운 원소를 추가해서 더
큰 해를 만드는 것이다.

\subsection{Initialization}\label{initialization}

첫 번째 반복이 시작될 때 \texttt{j\ =\ 2}이다.

이때 \texttt{A{[}1..j-1{]}\ =\ A{[}1{]}}인데, 원소가 하나뿐인 배열은
항상 정렬되어 있다.

따라서 loop invariant는 처음부터 성립한다.

\subsection{Maintenance}\label{maintenance}

반복이 시작될 때 \texttt{A{[}1..j-1{]}}이 정렬되어 있다고 가정하자.

알고리즘은 \texttt{key\ =\ A{[}j{]}}를 이 정렬된 구간 안의 적절한 위치에
삽입한다. 삽입 과정에서는 \texttt{key}보다 큰 원소들만 오른쪽으로
밀린다. 그 후 \texttt{key}가 정확한 위치에 놓인다.

그러면 결과적으로 \texttt{A{[}1..j{]}}가 정렬된다.

즉, 한 단계가 끝난 뒤에도 다음 단계에서 필요한 정렬 조건이 계속
유지된다.

\subsection{Termination}\label{termination}

반복문이 끝나면 \texttt{j\ =\ A.length\ +\ 1}이다.

그때 loop invariant에 의해 \texttt{A{[}1..A.length{]}}가 정렬되어
있으므로, 전체 배열이 올바르게 정렬된 것이다.

이로써 Insertion Sort의 correctness가 증명된다.

\section{왜 이것이 Incremental
Algorithm인가}\label{uxc65c-uxc774uxac83uxc774-incremental-algorithmuxc778uxac00}

이제 다시 핵심 질문으로 돌아가 보자.

왜 Insertion Sort를 incremental algorithm이라고 부르는가?

그 이유는 다음 문장으로 정리할 수 있다.

\begin{tcolorbox}[enhanced jigsaw, arc=.35mm, bottomrule=.15mm, bottomtitle=1mm, breakable, colback=white, colbacktitle=quarto-callout-tip-color!10!white, colframe=quarto-callout-tip-color-frame, coltitle=black, left=2mm, leftrule=.75mm, opacityback=0, opacitybacktitle=0.6, rightrule=.15mm, title=\textcolor{quarto-callout-tip-color}{\faLightbulb}\hspace{0.5em}{한 문장 요약}, titlerule=0mm, toprule=.15mm, toptitle=1mm]

Insertion Sort는 \textbf{이미 정렬된 부분문제의 해를 유지한 채, 새 원소
하나를 삽입하여 더 큰 부분문제의 해를 만드는 알고리즘}이다.\\

\end{tcolorbox}

좀 더 풀어서 말하면 다음과 같다.

\begin{itemize}
\item
  \texttt{A{[}1{]}}은 이미 정렬되어 있다.
\item
  \texttt{A{[}2{]}}를 삽입해서 \texttt{A{[}1..2{]}}를 만든다.
\item
  \texttt{A{[}3{]}}를 삽입해서 \texttt{A{[}1..3{]}}을 만든다.
\item
  \ldots{}
\item
  \texttt{A{[}n{]}}을 삽입해서 \texttt{A{[}1..n{]}}을 만든다.
\end{itemize}

즉, 정렬 결과를 매 단계 새로 계산하는 것이 아니라, \textbf{이전 단계의
결과를 재사용하여 한 단계씩 확장한다.}

이것이 incremental algorithm의 본질이다.

\section{시간 복잡도}\label{uxc2dcuxac04-uxbcf5uxc7a1uxb3c4}

Insertion Sort의 시간 복잡도는 입력 배열의 초기 상태에 따라 달라진다.

\subsection{Best Case}\label{best-case}

배열이 이미 정렬되어 있는 경우를 생각해 보자.

이 경우 각 단계에서 \texttt{A{[}i{]}\ \textgreater{}\ key} 조건이 곧바로
거짓이 되므로, \texttt{while} 루프는 거의 실행되지 않는다.

따라서 전체 수행 시간은 선형 시간에 가깝고,

\[
O(n)
\]

이 된다.

\subsection{Worst Case}\label{worst-case}

배열이 완전히 역순으로 정렬되어 있는 경우를 생각해 보자.

예를 들어

\begin{verbatim}
[6, 5, 4, 3, 2, 1]
\end{verbatim}

과 같은 경우, 각 단계에서 새 원소는 가장 앞쪽으로 들어가야 한다. 따라서
이전 원소들을 전부 오른쪽으로 밀어야 한다.

이때 비교 및 이동 횟수는 대략

\(1 + 2 + 3 + \cdots + (n−1)\)

이 되며, 이는

\[
O(n^2)
\]

에 해당한다.

\subsection{직관적 해석}\label{uxc9c1uxad00uxc801-uxd574uxc11d}

Insertion Sort는 각 단계에서 정렬된 구간 안에 새 원소를 삽입한다. 그런데
삽입 위치를 찾기 위해 최악의 경우 정렬된 구간 전체를 왼쪽 끝까지
확인해야 한다.

즉,

\begin{itemize}
\tightlist
\item
  첫 번째 삽입은 최대 1번 비교,
\item
  두 번째 삽입은 최대 2번 비교,
\item
  세 번째 삽입은 최대 3번 비교,
\item
  \ldots{}
\item
  (n-1)번째 삽입은 최대 (n-1)번 비교
\end{itemize}

처럼 진행되므로 전체 비용이 이차적으로 커질 수 있다.

\section{다른 설계 관점과의
비교}\label{uxb2e4uxb978-uxc124uxacc4-uxad00uxc810uxacfcuxc758-uxbe44uxad50}

Insertion Sort의 incremental 성격을 더 분명하게 보기 위해, 다른 알고리즘
설계 방식과 비교해 볼 수 있다.

\subsection{Incremental vs
Divide-and-Conquer}\label{incremental-vs-divide-and-conquer}

\subsubsection{Incremental}\label{incremental}

\begin{itemize}
\tightlist
\item
  이미 구해 놓은 작은 해를 바탕으로 다음 해를 확장한다.
\item
  부분 정답을 계속 유지한다.
\item
  예: Insertion Sort
\end{itemize}

\subsubsection{Divide-and-Conquer}\label{divide-and-conquer}

\begin{itemize}
\tightlist
\item
  문제를 더 작은 여러 부분문제로 나눈다.
\item
  각각을 독립적으로 해결한다.
\item
  마지막에 그 해들을 합친다.
\item
  예: Merge Sort
\end{itemize}

즉, Insertion Sort는 문제를 분할해서 푸는 것이 아니라,\\
\textbf{부분 정답을 유지하면서 점차 크게 만들어 가는 방식}이다.

이 차이를 이해하면 왜 Insertion Sort를 incremental algorithm이라고
부르는지 더 선명하게 보인다.

\section{Pseudo code}\label{pseudo-code}

아래는 Insertion Sort의 pseudo code이다.

\[
\begin{aligned}
&\text{INSERTION-SORT}(A) \\
&1\quad \mathbf{for}\ j = 2\ \mathbf{to}\ A.\mathrm{length} \\
&2\qquad \mathtt{key} = A[j] \\
&3\qquad \text{// Insert } A[j] \text{ into the sorted sequence } A[1 \ldots j-1] \\
&4\qquad i = j - 1 \\
&5\qquad \mathbf{while}\ i > 0 \ \mathbf{and}\ A[i] > \mathtt{key} \\
&6\qquad\qquad A[i+1] = A[i] \\
&7\qquad\qquad i = i - 1 \\
&8\qquad A[i+1] = \mathtt{key}
\end{aligned}
\]

\section{Conclusion}\label{conclusion}

Insertion Sort는 incremental algorithm의 대표적인 예시로, 각 단계에서
이미 정렬된 부분배열에 새 원소를 삽입하여 점진적으로 전체 배열을
정렬하는 방식으로 동작한다.\\
Loop invariant를 통해 알고리즘의 correctness를 증명할 수 있으며, 입력
배열의 초기 상태에 따라 시간 복잡도가 달라진다.\이러한 특성 때문에
Insertion Sort는 작은 배열이나 거의 정렬된 배열에 대해 효율적이며,
incremental algorithm의 개념을 이해하는 데 좋은 사례가 된다.




\end{document}
